\section{Discussion}
\paragraph{Conclusion.} We present Splat-SAP, a feed-forward approach for novel view synthesis of human-centered scenes.
%
In particular, we employ a 3D foundation model and utilize iterative affinity learning to reconstruct scale-aware point maps as a coarse geometry without 3D supervision.
%
We further leverage geometric constraints to refine the initial geometry, on which we build a Gaussian plane for rendering.
%
The full coarse-to-fine pipeline can be trained with only rendering loss by using multi-view image datasets.
%
Our method achieves superior rendering results with respect to baseline methods, especially in the case of sparse input views.

\paragraph{Limitation.} We notice the floating artifacts in Fig.~\ref{fig:abla_full}(e).
%
This is because MASt3R might predict some floating points on the boundary between foreground and background.
%
Since such regions are observed by only one of two input views, they can not be corrected by our refinement module.
%
We believe that incorporating monocular prior~\cite{xu2024depthsplat} would alleviate this problem.

\paragraph{Acknowledgement.} This paper is supported by National
Key R\&D Program of China (2022YFF0902200) and NSFC
project (Grants 62125107 and 62301298).